Redefining the International System of Units (SI) second requires
independent frequency-ratio measurements between different types of
optical clocks in different laboratories, with an overall measurement
uncertainty below \refThreshold{} and consistent results. However, the
self-assessed uncertainties of these clocks have not yet been validated
by independent frequency-ratio comparisons. Here we compare a ytterbium
optical lattice clock in Wuhan with a strontium optical lattice clock
in Shanghai---approximately $\siteKm$\,km away---using a
phase-stabilized fiber link spanning $\fibreKm$\,km. The Yb/Sr
frequency ratio is determined as \Rheadline(\RheadlineUnc), with a
Bayesian statistical uncertainty of \uStatShort$\times10^{-18}$ and a
total uncertainty of \uTotalE$\times10^{-18}$. The value differs from
that of NIST--JILA~\refNistV(\refNistU) by \dVsNistE{}, consistent
within the margin of statistical error. To our knowledge, this is the
first intercity, cross-species frequency-ratio measurement to meet the
uncertainty requirements outlined in the roadmap for redefining the SI
second, and together with the NIST--JILA result it satisfies the
mandatory criteria (I.2.b) for the redefinition, confirming the
validity of uncertainty assessments for both ytterbium and strontium
optical clocks.
